
\frontslide{L'arbre B}

%------------------------------------
%------------------------------------------------------------
\begin{frame}{Arbre-B}

Aboutissement des structures d'index basées sur \textbf{l'ordre}
des données

\vfill

\begin{redbullet}
  \item c'est un arbre équilibré
   \item chaque n{\oe}ud est un index local
  \item il se réorganise dynamiquement
\end{redbullet}

\vfill

Utilisé universellement\,!

\vfill
\begin{blocgauche}
Comparable aux séquentiel indexé, mais évite d'avoir
à maintenir des fichiers triés.
\end{blocgauche}
\end{frame}
%------------------------------------------------------------
\begin{frame}{Exemple d'un arbre B}

\figSlideWithSize{arbreBAnnees}{12cm}

\end{frame}
%------------------------------------------------------------


\begin{frame}{N{\oe}ud feuille d'un arbre B}

Un n{\oe}ud feuille est un \alert{index dense local}, contenant des entrées d'index.

\medskip

\figSlideWithSize{arbreBFeuille}{9cm}

\vfill

\begin{blocgauche}
Chaque entrée référence un (ou plusieurs) enregistrements du fichier de données:
celui (ceux) ayant la même valeur de clé que l'entrée.
\end{blocgauche}

\end{frame}

\begin{frame}{N{\oe}ud interne d'un arbre B}


Un n{\oe}ud interne est un \alert{index non dense local}, les enregistrements
servant de clé, intercalés avec des pointeurs.

\medskip

\figSlideWithSize{arbreBInterne}{10cm}

\vfill
\begin{blocgauche}
Le sous-arbre référencé par $A_2$ contient tous
les enregistrements dont la clé est comprise
entre $C_1$ et $C_2$.
\end{blocgauche}
\end{frame}


\begin{frame}{Avec notre fichier de 1M de films}

Admettons qu'une 
entrée d'index occupe 12 octets, soit 8 octets pour l'adresse, et 4 pour la clé (l'année du film).
    
    \vfill
    
    Chaque bloc contient 4 096 octets. On place donc $\lfloor \frac{4096}{12} \rfloor = 341$ entrées
    (au maximum) dans un bloc.
   
   
\begin{blocgauche}
\begin{redbullet}
    
    \item  Il faut $\lceil \frac{1000000}{341} \rceil=2\,933$ blocs 
    pour le  niveau des feuilles. 
   
   \item  Le deuxième niveau est \alert{non dense}. Il comprend autant d'entrées que de blocs
    à indexer, soit 2933.  Il faut
    donc $\lceil \frac{2\,933}{341} \rceil = 9$ blocs (au mieux).
    
    \item Finalement, un troisième niveau, constitué d'un bloc
     avec 8 entrées suffit pour compléter l'index.
\end{redbullet}
\end{blocgauche}

\end{frame}

%------------------------------------------------------------
\begin{frame}{Capacité d'un arbre B}


\begin{redbullet}
   \item  avec un niveau d'index (la racine seulement)
             on peut donc référencer 341 films\,;

   \item  avec deux niveaux on indexe 341
          blocs de 341 films chacun, soit $341^2 = 1\,162\,816 $ films\,;

   \item  avec trois niveaux on indexe $341^3 = 39\,651\,821$ films\,;

  \item enfin avec quatre niveaux on indexe plus de 1 milliard de films.
\end{redbullet}
\vfill

\begin{blocgauche}
\begin{redbullet}
  \item La hauteur de l'arbre est \alert{logarithmique} par rapport au nombre d'enregistrements.
  \item Inversement, avec un arbre de hauteur $h$, on indexe un collection de taille
      \alert{exponentielle} en $h$.
\end{redbullet}
\end{blocgauche}
\end{frame}


\begin{frame}{Insertion dans un arbre B}

On construit l'index sur le titre des films. Situation initiale.

\vfill

\figSlideWithSize{arbreB-1}{8cm}

\end{frame}

\begin{frame}{La procédure d'éclatement}

\figSlideWithSize{arbreB-2}{12cm}

\vfill
\begin{blocgauche}
Quand un n{\oe}ud est plein, il faut effectuer un \alert{éclatement}.
\end{blocgauche}
\end{frame}

\begin{frame}{Les insertions continuent}

\figSlideWithSize{arbreB-3}{12cm}

\vfill
\begin{blocgauche}
Underground, puis Easy Rider.
\end{blocgauche}
\end{frame}



\begin{frame}{Ajout d'un niveau}

\figSlideWithSize{arbreB-4}{12cm}

\vfill
\begin{blocgauche}
L'éclatement de la racine entraîne l'ajout d'un niveau
\end{blocgauche}
\end{frame}

\begin{frame}{Les insertions continuent}

\figSlideWithSize{arbreB-5}{12cm}

\vfill
\begin{blocgauche}
Shining et Annie Hall
\end{blocgauche}
\end{frame}

\begin{frame}{Les insertions continuent}

\figSlideWithSize{arbreB-6}{12cm}

\vfill
\begin{blocgauche}
Après insertion de Jurassic Park, 
Manhattan et Metropolis
\end{blocgauche}
\end{frame}

%%%%%%%%%%
\begin{comment}
\begin{frame}
\frametitle{Insertion dans un arbre-B d'ordre $k$}

On recherche la feuille de l'arbre o\`u l'enregistrement doit prendre place et
on l'y insère. Si le bloc $p$ déborde (il contient $2k+1$ éléments)\,:
\begin{redbullet}

  \item On alloue un nouveau bloc $p'$.

  \item On place les $k$ premiers enregistrements (ordonnés selon la clé)
        dans $p$ et les $k$ derniers dans $p'$.

  \item On insère le $k+1^{e}$ enregistrement dans le père de $p$. Son
        pointeur gauche référence $p$, et son pointeur droit
        référence $p'$.

  \item Si le père déborde \`a son tour, on continue
        comme en 1.

\end{redbullet}

\end{frame}
\end{commennt}


%------------------------------------------------------------

\begin{frame}
\frametitle{Recherche dans un arbre-B}

Rechercher les enregistrements de clé $C$.
\vfill

À partir de la racine,  appliquer récursivement 
l'algorithme suivant\,:

\vfill

Soit $C_{1}, \ldots C_{n}$ les valeurs de clés du bloc courant.

\begin{redbullet}

  \item Chercher $C$ dans le bloc courant. Si $C$ y est, accéder aux
valeurs des autres champs, Fin.

  \item Sinon, Si $C<C_{1}$ (ou $C>C_{n}$), on continue la recherche avec 
        le n{\oe}ud référencé par $P_{1}$ (ou $P_{n+1}$).

  \item Sinon, il existe $i\in \lbrack1,k\lbrack$ tel que
        $C_{i}< C < C_{i+1}$, on continue avec le bloc
        référencé par le pointeur $P_{i+1}$.

\end{redbullet}


\end{frame}

%

\begin{comment}
\begin{frame}
\frametitle{Quelques mesures pour l'arbre-B}

Hauteur $h$ d'un arbre-B d'ordre $k$ contenant $n$ articles\,:

$$\log_{2k+1}(n+1) \le h \le \log_{k+1}(\frac{n+1}{2})$$

\begin{redbullet}
  \item $\log_{2k+1}(n+1)$ : meilleur des cas (arbre balancé)
  \item $\log_{k+1}(\frac{n+1}{2})$ : pire des cas (insertion dans la même
  branche)
\end{redbullet}

Exemple pour $k=100$\,:

\begin{redbullet}
  \item si $h=2$, $n\le8 \times 10^{6}$
  \item si $h=3$, $n\le1,6 \times 10^{9}$
\end{redbullet}

Les opérations d'accès co\^utent au maximum $h$ E/S (en moyenne
$\geq h-1$).

\end{frame}
\end{comment}


%------------------------------------------------------------
\begin{frame}[fragile]{Recherche par clé}

\begin{minted}[frame=leftline,
               baselinestretch=.8,
               fontsize=\footnotesize,
               framesep=2mm]{sql}
SELECT *
FROM  Film
WHERE titre = 'Manhattan'
\end{minted}

\begin{redbullet}
   \item on lit la racine de l'arbre\,: {\it Manhattan} étant situé dans
   l'ordre lexicographique entre {\it Easy Rider} et {\it Psychose},
   on doit suivre le chaînage situé entre ces deux titres\,;

  \item on lit le bloc intermédiaire, on descend à gauche;
  
  \item dans la feuille, on trouve l'entrée correspondant à Manhattan;

  \item il reste à lire l'enregistrement. 
\end{redbullet}

\end{frame}
%------------------------------------------------------------
\begin{frame}[fragile]{Recherche par intervalle}

\begin{minted}[frame=leftline,
               baselinestretch=.8,
               fontsize=\footnotesize,
               framesep=2mm]{sql}
select *
from  Film
where annee between 1960 AND 1975
\end{minted}

\begin{redbullet}
   \item On fait une recherche par clé pour l'année 60

  \item on parcourt les feuilles de l'arbre en suivant le chaînage,
            jusqu'à l'année 1975

  \item à chaque fois on lit l'enregistrement
\end{redbullet}

Attention, les accès aux fichiers peuvent coûter très cher.

\end{frame}

\begin{frame}{L'arbre sur les années}

\figSlideWithSize{arbreBAnnees}{12cm}
\end{frame}

%------------------------------------------------------------
\begin{frame}[fragile]{Recherche par préfixe}

Exemple\,:

\begin{minted}[frame=leftline,
               baselinestretch=.8,
               fontsize=\footnotesize,
               framesep=2mm]{sql}
SELECT *
FROM  Film
WHERE titre LIKE 'M\%'
\end{minted}

\vfill

Revient à une recherche par intervalle.\\

\vfill

\begin{minted}[frame=leftline,
               baselinestretch=.8,
               fontsize=\footnotesize,
               framesep=2mm]{sql}
SELECT *
FROM  Film
WHERE titre BETWEEN  'MAAAAAA...' AND 'MZZZZZZ...'
\end{minted}


\vfill

Contre-exemple\,:

\begin{minted}[frame=leftline,
               baselinestretch=.8,
               fontsize=\footnotesize,
               framesep=2mm]{sql}
SELECT *
FROM  Film
WHERE titre LIKE '\%e'
\end{minted}
\vfill

 Ici index inutilisable


\end{frame}



%------------------------------------------------------------
\begin{frame}[fragile]{Création d'un arbre B}

Sur la clé primaire.

\begin{minted}[frame=leftline,
               baselinestretch=.8,
               fontsize=\footnotesize,
               framesep=2mm]{sql}
   create table Film (titre varchar(30) not null,
                      ...,
                      primary key (titre)
                      );
\end{minted}
\vfill

Sur n'importe quel attribut.

\begin{minted}[frame=leftline,
               baselinestretch=.8,
               fontsize=\footnotesize,
               framesep=2mm]{sql}
   create index filmAnnee on Film (annee)
\end{minted}

\vfill

Le SGBD synchronise le contenu de la table et celui de l'index
\end{frame}
%------------------------------------------------------------
\begin{frame}{Résumé: efficacité de l'arbre B}


Il est (presque) parfait\,!

\begin{redbullet}
  \item On a très rarement besoin de plus de trois niveaux
  \item Le coût d'une recherche par clé est le nombre de niveaux, plus 1.

  \item Supporte les recherches par clé, par intervalle,
             par préfixe
  \item Dynamique
\end{redbullet}
On peut juste lui reprocher d'occuper de la place.

\end{frame}